% !TeX spellcheck = ru_RU
\documentclass[a4paper, 14pt, unknownkeysallowed]{extreport}
\usepackage[utf8]{inputenc}
\usepackage[T2A]{fontenc}
\usepackage[english,russian]{babel}
\usepackage{cmap}
\usepackage{enumitem}

\usepackage{csvsimple}
\usepackage{hyphenat} 
\usepackage{amstext, amsmath,amsfonts,amssymb,amsthm,mathtools} 

\usepackage{setspace}
\onehalfspacing 

\usepackage{geometry}
\geometry{left=30mm}
\geometry{right=10mm}
\geometry{top=10mm}
\geometry{bottom=20mm}

\usepackage{graphicx}
\graphicspath{{img/}} 

\usepackage{threeparttable}
\usepackage{bigdelim}

\usepackage{setspace}
\onehalfspacing % Полуторный интервал

\frenchspacing
\usepackage{indentfirst} % Красная строка

\usepackage{cases}

\usepackage[unicode,pdftex]{hyperref} % Ссылки в pdf
\hypersetup{hidelinks}

\usepackage{xcolor}
\usepackage{listings}
% Для листинга кода:
\lstset{%
	language=Matlab,   					% выбор языка для подсветки	
	basicstyle=\small\sffamily,			% размер и начертание шрифта для подсветки кода
	numbers=left,						% где поставить нумерацию строк (слева\справа)
	%numberstyle=,					% размер шрифта для номеров строк
	stepnumber=1,						% размер шага между двумя номерами строк
	numbersep=5pt,						% как далеко отстоят номера строк от подсвечиваемого кода
	frame=single,						% рисовать рамку вокруг кода
	tabsize=4,							% размер табуляции по умолчанию равен 4 пробелам
	captionpos=t,						% позиция заголовка вверху [t] или внизу [b]
	breaklines=true,					
	breakatwhitespace=true,				% переносить строки только если есть пробел
	escapeinside={\#*}{*)},				% если нужно добавить комментарии в коде
	backgroundcolor=\color{white},
}
\usepackage{caption}
\captionsetup{labelsep = endash}
\captionsetup[figure]{name = {Рисунок}, justification = centerlast}
\captionsetup[lstlisting]{justification = justified}
\captionsetup[table]{justification = justified}

\usepackage{titlesec}
\titleformat{\chapter}{\LARGE\bfseries}{\thechapter}{20pt}{\LARGE\bfseries}
\titleformat{\section}{\Large\bfseries}{\thesection}{20pt}{\Large\bfseries}

\renewcommand\labelitemi{\ --\ }
\renewcommand\labelenumi{\theenumi)}

\newcommand{\img}[3] {
	\begin{center}
		\begin{figure}[h!]
			\includegraphics[scale = #1]{img/#2}
			\caption{#3}
			\label{img:#2}
		\end{figure}
	\end{center}
}

\newcommand{\lst}[4]{
\lstinputlisting[language=Matlab, firstline=#1, lastline=#2, label=lst:#3,caption=#4]{../src/code.m}
}

\usepackage{siunitx,array,booktabs}
\usepackage{amsthm}
\usepackage{tabularx}
\newtheorem{theorem}{Theorem}
\newtheorem{definition}{Опредление}


\begin{document}
% !TeX spellcheck = ru_RU
\begin{titlepage}
	
	\newgeometry{pdftex, left=2cm, right=2cm, top=2.5cm, bottom=2.5cm}
	\fontsize{12pt}{12pt}\selectfont
	\noindent \begin{minipage}{0.15\textwidth}
		\includegraphics[width=\linewidth]{img/b_logo.jpg}
	\end{minipage}
	\noindent\begin{minipage}{0.9\textwidth}\centering
		\textbf{Министерство науки и высшего образования Российской Федерации}\\
		\textbf{Федеральное государственное бюджетное образовательное учреждение высшего образования}\\
		\textbf{«Московский государственный технический университет имени Н.~Э.~Баумана}\\
		\textbf{(национальный исследовательский университет)»}\\
		\textbf{(МГТУ им. Н.~Э.~Баумана)}
	\end{minipage}

	\noindent\rule{18cm}{3pt}
	\newline\newline
	\noindent ФАКУЛЬТЕТ $\underline{\text{«Информатика и системы управления»~~~~~~~~~~~~~~~~~~~~~~~~~~~~~~~~~~~~~~~~~~~~~~~~~~~~~~~}}$ \newline\newline
	\noindent КАФЕДРА $\underline{\text{«Программное обеспечение ЭВМ и информационные технологии»~~~~~~~~~~~~~~~~~~~~~~~}}$\newline\newline\newline\newline\newline\newline\newline
	
	
	\begin{center}
		\Large\textbf{Отчет по лабораторной работе № 4}
		
		\Large\textbf{по курсу <<Моделирование>>}
		
		\textbf{\newline}		
	\end{center}
	
	\noindent\textbf{Тема} $\underline{\text{~~Моделирование простейшей СМО~~~~~~~~~~~~~~~~~~~~~~~~~~~~~~~~~~~~~~~~~~~~~~~~~~~~~~~~~~~~~~~~~~~~~~~~}}$\newline\newline
	\noindent\textbf{Студент} $\underline{\text{~~Виноградов А. О.~~~~~~~~~~~~~~~~~~~~~~~~~~~~~~~~~~~~~~~~~~~~~~~~~~~~~~~~~~~~~~~~~~~~~~~~~~~~~~~~~~~~~~}}$\newline\newline
	\noindent\textbf{Группа} $\underline{\text{~~ИУ7-76Б~~~~~~~~~~~~~~~~~~~~~~~~~~~~~~~~~~~~~~~~~~~~~~~~~~~~~~~~~~~~~~~~~~~~~~~~~~~~~~~~~~~~~~~~~~~~~~~~~~~~}}$\newline\newline
	\noindent\textbf{Оценка (баллы)} $\underline{\textbf{~~~~~~~~~~~~~~~~~~~~~~~~~~~~~~~~~~~~~~~~~~~~~~~~~~~~~~~~~~~~~~~~~~~~~~~~~~~~~~~~~~~~~~~~}}$\newline\newline
	\noindent\textbf{Преподаватель} $\underline{\text{~~Рудаков И. В.~~~~~~~~~~~~~~~~~~~~~~~~~~~~~~~~~~~~~~~~~~~~~~~~~~~~~~~}}$\newline
	
	\begin{center}
		\vfill
		Москва~---~\the\year
		~г.
	\end{center}
	\restoregeometry
	
\end{titlepage}
\setcounter{page}{2}

\chapter{Теоретическая часть}
Случайный процесс называется марковским, если он обладает следующим свойством: для каждого момента $t_0$ вероятность любого состояния системы в будущем (при $t > t_0$) зависит только от ее состояния в настоящем (при $t = t_0$) и не зависит от того, когда и каким образом система пришла в это состояние.  

Вероятностью i-го состояния называется вероятность $p_i(t)$ того, что в момент t система будет находится в состоянии $S_i$. Также для любого времени t выполняется условие нормировки --- сумма вероятностей всех состоянии равна единице.

Стационарное распределение цепи Маркова --- распределение вероятности, которое не меняется с течением времени.

Предельная вероятность состояния показывает среднее время пребывания системы в этом состоянии. Для нахождения таких вероятностей используют уравнения Колмогорова.

Уравнение Колмогорова строится по следующему правилу: в левой части каждого уравнения стоит производная вероятности состояния, а правая часть содержит столько членов сколько стрелок связано с данным состоянием.

Если стрелка направлена из состояния, то соответствующий член имеет знак "$-$", в состояние --- "$+$". 

Каждый член равен произведению плотности вероятности перехода соответствующий данной стрелке, умноженной на вероятность того состояния, из которого исходит стрелка. 

\newpage
Получаем систему уравнений вида:

\begin{equation}
	\label{ass}
	\begin{cases}
		p_1'(t) =  -(\lambda_{11} + ... + \lambda_{1n})p_1 + \lambda_{21}*p_2+...+\lambda_{n1}*p_n\\
		p_2'(t) =  -(\lambda_{21} + ... + \lambda_{2n})p_2 + \lambda_{12}*p_1+...+\lambda_{n2}*p_n\\
		...\\
		p_n'(t) =  -(\lambda_{n1}+ ... + \lambda_nn)p_n + \lambda_{1n}*p_1+...+\lambda_{nn}*p_{n}.
	\end{cases}
\end{equation} 

В стационарном случае ($p_i'(t) = 0, i = \overline{1, n}$) левые системы уравнений принимают значение 0. Получаем СЛАУ, решение которой дает нам искомые предельные вероятности состояний системы. 

Интегрирование системы (\ref{ass}) дает искомые вероятности системы как функции времени. 
График зависимости вероятности пребывания системы в состоянии $p_i$ колеблется, а затем в определенный момент $t_i!$ принимает стабильное значение $p_i!$. Это время $t_i!$ до стабилизации вычисляется в реализованной программе для каждого $i = \overline{1, n}$.

\chapter{Результат}
На рисунке \ref{img:1} приведен пример работы программы.
\img{0.8}{1}{Пример работы программы}

В листингах \ref{lst:1} -- \ref{lst:2} приведена реализация вычислений коэффициентов уравнений Колмогорова, вычисления предельных вероятностей и времени до стабилизации.
\lst{9}{13}{func.py}{1}{Реализация программных вычислений}
\lst{14}{54}{func.py}{2}{Реализация программных вычислений (продолжение)}

\end{document}